\documentclass{article}
\usepackage[utf8]{inputenc}
\usepackage[spanish]{babel}
\usepackage{amsmath}
\usepackage{enumitem}
\usepackage{circuitikz}
\usepackage{float}

% Solución para el conflicto entre circuitikz y babel en español
\usetikzlibrary{babel}

\title{PRÁCTICA 19 \\ LEY DE TENSIÓN DE KIRCHHOFF (PARA UN GENERADOR)}
\date{}

\begin{document}

\maketitle

\section*{FINALIDADES}
\begin{enumerate}
    \item Determinar por análisis o cálculo la relación entre la suma de las caídas de tensión entre los extremos de resistencias conectadas en serie y la tensión aplicada.
    \item Comprobar experimentalmente la relación determinada en 1.
\end{enumerate}

\section*{INFORMACIÓN PRELIMINAR}
La solución de circuitos eléctricos complicados resulta fácil aplicando las leyes de Kirchhoff. Estas leyes fueron formuladas y publicadas por el físico Gustav Robert Kirchhoff (1824-1887), y establecieron la base del análisis moderno de redes. Son aplicables a otros circuitos que contienen una o más fuentes de tensión. En esta práctica trataremos de los que sólo tienen una fuente de tensión.

\begin{figure}[H]
    \centering
    \begin{circuitikz}[american]
        \draw (0,0)
        to[battery, l=$V$, invert] (0,3)
        to[short, i>^={\color{red}$I_T$}] (2,3)
        to[R, l_=$R_1$, v^=$V_1$] (4,3)
        to[R, l_=$R_2$, v^=$V_2$] (6,3)
        to[R, l_=$R_3$, v^=$V_3$] (8,3)
        to[R, l_=$R_4$, v^=$V_4$] (10,3)
        to[short] (10,0)
        to[short] (0,0);
    \end{circuitikz}
    \caption{Circuito serie con cuatro resistencias y una fuente de tensión}
    \label{fig:circuito}
\end{figure}

Sustituyendo el valor equivalente \( R_T \) de la ecuación (19-3) en la ecuación (19-2) se tiene
\[
V = I_T(R_1 + R_2 + R_3 + R_4)
\]
\[
V = I_T \times R_1 + I_T \times R_2 + I_T \times R_3 + I_T \times R_4
\tag{19-4}
\]

Una de las características de un circuito serie es que la corriente es la misma en cualquier punto del circuito. \( I_T \) es pues la corriente total en \( R_1, R_2, R_3 \) y \( R_4 \). De esto se deduce que
\begin{align*}
I_T \times R_1 &= V_1, \text{ caída de tensión entre los extremos de } R_1 \\
I_T \times R_2 &= V_2, \text{ caída de tensión entre los extremos de } R_2 \\
I_T \times R_3 &= V_3, \text{ caída de tensión entre los extremos de } R_3 \\
\text{e } I_T \times R_4 &= V_4, \text{ caída de tensión entre los extremos de } R_4
\end{align*}

Ahora se puede escribir la ecuación (19-4) como sigue:
\[
V = V_1 + V_2 + V_3 + V_4
\tag{19-5}
\]

La ecuación (19-5) es la expresión matemática de la ley de tensión de Kirchhoff para resistencias conectadas en serie en un circuito cerrado. Es evidente que se puede generalizar la ec. (19-5) para circuitos que contienen una o más resistencias conectadas en serie en

\end{document}